\documentclass[11pt,a4paper]{article}
% Compile with XeLaTeX. Fonts reproduce the 2022 original: TeX Gyre Pagella text;
% Euler math letters (Latin and Greek, eulervm's zeur); newpx symbols and
% extension fonts; Latin Modern math operators and digits; bbold blackboard bold;
% Euler fraktur. The load order eulervm -> newpxmath -> Euler letters matches the
% set of fonts embedded in the original PDF.
\usepackage{amsmath,amssymb,amsthm,mathtools}
\usepackage{eulervm}
\usepackage{newpxmath}
\DeclareSymbolFont{letters}{U}{zeur}{m}{n}
\SetSymbolFont{letters}{bold}{U}{zeur}{b}{n}
% eulervm routes these to its own fonts; newpxmath does not re-declare them.
\DeclareMathSymbol{:}{\mathrel}{operators}{"3A}
\DeclareMathSymbol{\infty}{\mathord}{symbols}{"31}
\DeclareMathAlphabet{\mathfrak}{U}{euf}{m}{n}
\SetMathAlphabet{\mathfrak}{bold}{U}{euf}{b}{n}
\usepackage{bbold}
\usepackage[no-math]{fontspec}
\setmainfont{TeX Gyre Pagella}
\usepackage[margin=2.5cm]{geometry}
\usepackage{enumitem,booktabs,array}
\usepackage{tikz}
\usepackage{microtype}
\usepackage{xurl}
\usepackage[hidelinks]{hyperref}

\numberwithin{equation}{section}
\theoremstyle{plain}
\newtheorem{theorem}{Theorem}[section]
\newtheorem{lemma}[theorem]{Lemma}
\newtheorem{proposition}[theorem]{Proposition}
\newtheorem{corollary}[theorem]{Corollary}
\newtheorem*{theoremA}{Theorem A}
\newtheorem*{theoremB}{Theorem B}
\theoremstyle{definition}
\newtheorem{definition}[theorem]{Definition}
\newtheorem{construction}[theorem]{Construction}
\newtheorem{example}[theorem]{Example}
\newtheorem{question}{Question}
\theoremstyle{remark}
\newtheorem{remark}[theorem]{Remark}

\newcommand{\Q}{\mathbb{Q}}
\newcommand{\R}{\mathbb{R}}
\newcommand{\C}{\mathbb{C}}
\newcommand{\Z}{\mathbb{Z}}
\newcommand{\HT}{H^{*}_{T}}
\newcommand{\cL}{\mathcal{L}}
\newcommand{\cI}{\mathcal{I}}
\DeclareMathOperator{\Fl}{Fl}
\DeclareMathOperator{\Hess}{Hess}
\DeclareMathOperator{\Gr}{Gr}
\DeclareMathOperator{\ord}{ord}
\DeclareMathOperator{\rk}{rk}
\DeclareMathOperator{\cl}{cl}
\DeclareMathOperator{\Flats}{Flats}
\DeclareMathOperator{\Cl}{Cl}
\DeclareMathOperator{\GrH}{Gr}
\DeclareMathOperator{\St}{St}
\DeclareMathOperator{\cone}{cone}
\DeclareMathOperator{\trop}{trop}
\DeclareMathOperator{\Sym}{Sym}
\DeclareMathOperator{\diag}{diag}
\DeclareMathOperator{\Hom}{Hom}
\DeclareMathOperator{\OS}{OS}

\emergencystretch=3em
\hyphenation{Ay-zen-berg Che-re-pa-nov Bu-chs-ta-ber So-lo-ma-din equi-vari-antly}

\newtheorem*{theoremC}{Theorem C}
\begin{document}

\title{Face posets, local weight matroids and Betti numbers of torus actions}
\author{David A.~Mikheev}
\date{Draft, \today}
\maketitle

\begin{abstract}
Let a compact torus act on a closed manifold with isolated fixed points. Its face poset $S(X)$ is covered by the lattices of flats of the linear matroids of tangent weights at the fixed points. We study how much of the cohomology these combinatorial data see. First, the local equivariant cohomology $\HT(X,X\setminus pt)$, embedded in $\HT(pt)$ by restriction, is the principal ideal generated by the product of the tangent weights; it therefore determines the weights up to scalars, with multiplicities, and hence the local lattice $S(X)_{\geqslant pt}$. Second, the local weight matroids do not determine the equivariant cohomology: two connected regular semisimple Hessenberg varieties in $\Fl(\C^5)$, with actions of complexity one, have isomorphic weight matroids at all fixed points but different Betti numbers. Third, the global face poset does determine the Betti numbers for every regular semisimple Hessenberg variety of type~A, full or partial, and for the isospectral manifolds of Ayzenberg and Buchstaber. The reason is rigidity: the cluster-permutohedron $\Cl_\Gamma$ determines the graph $\Gamma$ up to isomorphism, which we prove using Whitney's line graph theorem.
\end{abstract}

%======================================================================
\section{Introduction}
%======================================================================

Throughout, a \emph{torus action} is an effective action of a compact torus $T=T^k$ on a closed smooth manifold $X=X^{2n}$ with a finite nonempty set of fixed points. Toric topology studies such actions of complexity $n-k$ zero by combinatorial means~\cite{BP15}. Buchstaber and Terzić extended these methods to positive complexity~\cite{BT19a,BT19b}, and Ayzenberg, Masuda and their coauthors made the \emph{face poset} $S(X)$ of the action the central combinatorial object~\cite{AC22,AMS23,AB24}. Its elements are the closures of the components of the orbit type strata of $X/T$ that contain a fixed point, ordered by inclusion. By a theorem of Ayzenberg and Cherepanov~\cite{AC22}, $S(X)$ is \emph{locally geometric}: for a fixed point $p$, the faces containing $p$ form the lattice of flats of the linear matroid of the tangent weights at~$p$.

This paper asks what these combinatorial data know about the cohomology of $X$. We prove three results, one local and two global.

\begin{theoremA}[Theorem~\ref{thm:main}]
Let $pt$ be a fixed point with tangent weights $\alpha_1,\dots,\alpha_n\in H^2(BT;\Q)$. Restriction identifies $\HT(X,X\setminus pt)$ with the principal ideal $(\alpha_1\cdots\alpha_n)$ of $\HT(pt)=H^*(BT;\Q)$. This ideal determines the multiset of weights up to nonzero scalars, hence the linear matroid of weights at $pt$ with its parallel classes, and the lattice $S(X)_{\geqslant pt}$.
\end{theoremA}

The point of Theorem~A is that no independence assumption on the weights is needed. The familiar GKM argument computes the ideal as an intersection $\bigcap_j(\alpha_j)$, which equals the product only when the weights are pairwise independent and loses multiplicities otherwise. We use the Thom isomorphism and unique factorization instead.

The local data are far from determining the global cohomology. The Chow ring~\cite{AHK18} and the Orlik--Solomon algebra~\cite{OT92} of a matroid depend only on the matroid, so the following shows that neither of them, taken at all fixed points, determines $\HT(X)$.

\begin{theoremB}[Theorem~\ref{thm:counterexample}]
Let $M=\diag(1,2,3,4,5)$. The regular semisimple Hessenberg varieties $X_1=\Hess(M,(2,3,5,5,5))$ and $X_2=\Hess(M,(2,4,4,5,5))$ in $\Fl(\C^5)$ are connected GKM manifolds with the same $120$ fixed points and torus actions of complexity one. At every fixed point of either variety the weight matroid is $U_{2,3}\oplus U_{1,1}\oplus U_{1,1}$. Their Poincaré polynomials are $1+17q+42q^2+42q^3+17q^4+q^5$ and $1+12q+47q^2+47q^3+12q^4+q^5$, where $q$ has degree~$2$.
\end{theoremB}

The global face posets of $X_1$ and $X_2$ are different. This leaves the question, raised in~\cite{Mik22}, whether the face poset determines the Betti numbers of an equivariantly formal action. In complexity zero it does: for a quasitoric manifold the Betti numbers are the $h$-numbers of the orbit polytope~\cite{BP15}. Our third result answers the question for the main family of examples in positive complexity.

\begin{theoremC}[Theorems~\ref{thm:cl-rigid} and~\ref{thm:q1-partial}, Corollary~\ref{cor:q1-hess}]
\begin{enumerate}[nosep,label=(\arabic*)]
\item For simple graphs $\Gamma,\Gamma'$, an isomorphism of cluster-permutohedra $\Cl_\Gamma\cong\Cl_{\Gamma'}$ implies $\Gamma\cong\Gamma'$.
\item Let $X$ and $X'$ each be an isospectral manifold $M_{\Gamma,\lambda}$ or a regular semisimple Hessenberg variety $\Hess(M,h)$ of type~A, with $\Gamma$, resp.\ $\Gamma_h$, connected. If $S(X)\cong S(X')$, then $X$ and $X'$ have the same Betti numbers and, up to an automorphism of the torus, isomorphic equivariant cohomology.
\item The same conclusion about Betti numbers holds for partial regular semisimple Hessenberg varieties, including all partial flag varieties.
\end{enumerate}
\end{theoremC}

The cluster-permutohedron $\Cl_\Gamma$ is the face poset of the isospectral manifold $M_{\Gamma,\lambda}$, and through the twin construction of Ayzenberg and Buchstaber also of $\Hess(M,h)$ when $\Gamma=\Gamma_h$~\cite{AB21,AB24}. Part~(1) is purely combinatorial. Counting the minimal elements below joins of pairs of atoms recovers the line graph of $\Gamma$, and Whitney's theorem~\cite{Whi32} recovers $\Gamma$. For partial flags we need Jung's strengthening of Whitney's theorem~\cite{Jun66}, a fibration over the partial Hessenberg variety, and a finite computation for components with at most four vertices.

Theorem~C does not give a formula for the Betti numbers in terms of the poset. It shows that in these families the face poset remembers the whole graph, and the graph remembers the manifold. A counterexample to the general question, if one exists, must therefore come from outside type~A.

\paragraph{Computations and formalization.}
All computations cited in this paper were done with independent programs, which are available with their output at \url{https://github.com/davidmiheev/matroid_toric}. The Betti numbers in Theorem~B were computed both from the Białynicki-Birula formula and, independently, from the GKM presentation of $\HT$ in Macaulay2~\cite{M2}. Part~(1) of Theorem~C was confirmed by computing exact poset isomorphism classes, with nauty's Traces~\cite{MP14}, for all $208$ graphs with at most $6$ vertices. Before the proof, Theorem~C(2) and~(3) were verified by direct computation for all full and partial Hessenberg varieties with $n\leqslant7$ (714 and 2\,940 cases). The algebraic step of Theorem~A, Lemma~\ref{lem:ufd}(2), is formalized in Lean~4 with Mathlib in the same repository, together with the lattice of flats of a matroid and Whitney's Boolean expansion of its characteristic polynomial.

\paragraph{Conventions.}
Cohomology has rational coefficients. We identify $H^*(BT)=\Sym(\mathfrak t^*_\Q)\cong\Q[t_1,\dots,t_k]$ with $\deg t_i=2$ and $H^2(BT)=\Hom(T,S^1)\otimes\Q$. For a polynomial in $q$ we write $q$ for the class of degree~$2$, so $P_X(q)=\sum_j\beta_{2j}(X)q^j$ when $H^{\mathrm{odd}}(X)=0$.

%======================================================================
\section{Preliminaries}\label{sec:prelim}
%======================================================================

\subsection{Weights, face posets and GKM manifolds}
Let $p\in X^T$. The tangent representation decomposes as $T_pX\cong\bigoplus_{j=1}^nV(\alpha_{p,j})$, where $V(\alpha)=\C$ with $t\cdot z=\alpha(t)z$. Without an invariant complex structure the weights $\alpha_{p,j}\in\Hom(T,S^1)$ are defined up to sign, which does not affect linear dependence. The point $p$ is isolated if and only if no weight vanishes. For an oriented $T$-vector bundle, the equivariant Euler class of $T_pX$ at a fixed point is $e_T(T_pX)=\pm\prod_j\alpha_{p,j}\in H^{2n}(BT)$.

Let $X_j\subset X$ be the union of the orbits of dimension at most $j$. A \emph{face} of rank $j$ is the closure of a connected component of $X_j/T\setminus X_{j-1}/T$ that contains a fixed point, and $S(X)$ is the set of faces ordered by inclusion.

\begin{theorem}[Ayzenberg--Cherepanov~{\cite[Thm.~1]{AC22}}]\label{thm:AC}
For every $p\in X^T$, the poset $S(X)_{\geqslant p}$ of faces containing $p$ is isomorphic to the lattice of flats of the linear matroid of the weights $\alpha_{p,1},\dots,\alpha_{p,n}\in\Hom(T,S^1)\otimes\Q$. In particular $S(X)$ is graded, and each $S(X)_{\geqslant p}$ is a geometric lattice.
\end{theorem}

The action is \emph{equivariantly formal} if the Serre spectral sequence of $X_T\to BT$ degenerates; for $X^T$ finite and nonempty this holds if and only if $H^{\mathrm{odd}}(X)=0$, and then $\HT(X)\cong H^*(X)\otimes H^*(BT)$ as graded $H^*(BT)$-modules. A \emph{GKM manifold} is an equivariantly formal $X$ whose weights at each fixed point are pairwise linearly independent. Its fixed points and invariant $2$-spheres form the \emph{GKM graph}, whose edge $pq$ is labelled by the weight $\alpha_{pq}$ along the sphere, and by the GKM theorem~\cite{GKM98,CS74}
\[
\HT(X)\cong\bigl\{(f_p)\in\textstyle\bigoplus_{p\in X^T}H^*(BT)\bigm| f_p-f_q\in(\alpha_{pq})\text{ for every edge }pq\bigr\}.
\]

\subsection{Hessenberg varieties}
A \emph{Hessenberg function} is a non-decreasing map $h\colon[n]\to[n]$ with $h(j)\geqslant j$. For $M=\diag(\lambda_1,\dots,\lambda_n)$ with pairwise distinct real $\lambda_i$, the \emph{regular semisimple Hessenberg variety} is
\[
\Hess(M,h)=\{V_\bullet\in\Fl(\C^n)\mid MV_i\subset V_{h(i)}\text{ for all }i\}.
\]
It is a smooth projective variety of complex dimension $\sum_j(h(j)-j)$ with $H^{\mathrm{odd}}=0$~\cite{DPS92,AH20}, invariant under the diagonal torus $T=T^n$. The fixed points are the $n!$ coordinate flags, indexed by $w\in S_n$, and the tangent weights at $w$ are $t_{w(i)}-t_{w(j)}$ for $j<i\leqslant h(j)$~\cite{DPS92,Tym08}. Let $\Gamma_h$ be the graph on $[n]$ with edges $\{j,i\}$, $j<i\leqslant h(j)$. Up to isomorphism, the graphs $\Gamma_h$ are exactly the indifference (unit interval) graphs~\cite{AB21}. By~\cite{DPS92},
\begin{equation}\label{eq:dps}
P_{\Hess(M,h)}(q)=\sum_{w\in S_n}q^{\ell_h(w)},\qquad \ell_h(w)=\#\{(j,i)\mid j<i\leqslant h(j),\ w(j)>w(i)\}.
\end{equation}

Let $\bar k=(k_1,\dots,k_s)$ be a composition of $n$, $K_b=k_1+\dots+k_b$, and $h$ a Hessenberg function on $[s]$. The \emph{partial Hessenberg variety} $\Hess_{\bar k}(M,h)$ is the set of partial flags $V_{K_1}\subset\dots\subset V_{K_s}=\C^n$ with $MV_{K_a}\subset V_{K_{h(a)}}$ for all~$a$. Its fixed points are the ordered set partitions $J=(J_1,\dots,J_s)$ of $[n]$ with $|J_a|=k_a$, and its tangent weights at $J$ are $t_x-t_y$ for $x\in J_a$, $y\in J_b$, $a<b\leqslant h(a)$~\cite[§5]{AB24}. For a generic one-parameter subgroup, the Białynicki-Birula decomposition gives
\begin{equation}\label{eq:bb}
P_{\Hess_{\bar k}(M,h)}(q)=\sum_Jq^{\operatorname{ind}(J)},\qquad\operatorname{ind}(J)=\#\{(x,y)\mid x\in J_a,\ y\in J_b,\ a<b\leqslant h(a),\ x>y\},
\end{equation}
which reduces to~\eqref{eq:dps} for $\bar k=(1,\dots,1)$. The \emph{inflation} $\Gamma_h(\bar k)$ is the graph obtained from $\Gamma_h$ by replacing each vertex $a$ by $k_a$ pairwise non-adjacent copies, copies of $a$ and $b$ being adjacent iff $ab$ is an edge of $\Gamma_h$. At the fixed point $J$, identifying $x\in J_a$ with a copy of~$a$, the tangent weights are $t_x-t_y$ over the edges $xy$ of $\Gamma_h(\bar k)$, so the local weight matroid is the graphic matroid $M(\Gamma_h(\bar k))$.

\subsection{Isospectral manifolds and cluster-permutohedra}
Let $\Gamma$ be a connected graph on $[n]$ and $\lambda$ a generic spectrum. The \emph{isospectral manifold} $M_{\Gamma,\lambda}$ is the space of Hermitian matrices $A$ with spectrum $\lambda$ and $a_{ij}=0$ whenever $\{i,j\}\notin E(\Gamma)$, with $T^n$ acting by conjugation~\cite{AB24}. For $\Gamma=K_n$ it is $\Fl(\C^n)$.

\begin{definition}[{\cite[Def.~2.14]{AB24}}]\label{def:cluster}
Let $\Gamma$ be a simple graph on $[n]$. A \emph{clustering} is a partition $C=\{V_1,\dots,V_m\}$ of $[n]$ into blocks inducing connected subgraphs; ordered by refinement, the clusterings form the lattice of flats $\Flats(\Gamma)$ of the graphic matroid. An \emph{assignment} for $C$ is a family of label sets $L_1,\dots,L_m$ partitioning $[n]$ with $|L_i|=|V_i|$. The \emph{cluster-permutohedron} $\Cl_\Gamma$ is the set of pairs (clustering, assignment), where $(C',A')\leqslant(C,A)$ if $C'$ refines $C$ and the label set of each block of $C$ is the union of the label sets of the blocks of $C'$ it contains. It is graded by $\rk(C,A)=n-|C|$.
\end{definition}

\begin{theorem}[Ayzenberg--Buchstaber~{\cite[Thm.~1, Cor.~2.43]{AB24}}, \cite{AB21}]\label{thm:AB}
For connected $\Gamma$ and generic $\lambda$, $S(M_{\Gamma,\lambda})\cong\Cl_\Gamma$. If $\Gamma_h$ is connected, then $\Hess(M,h)$ and $M_{\Gamma_h,\lambda}$ are twins: they have homeomorphic orbit spaces by a homeomorphism preserving orbit dimensions, isomorphic equivariant cohomology and equal Betti numbers. In particular $S(\Hess(M,h))\cong\Cl_{\Gamma_h}$. For partial Hessenberg varieties, $S(\Hess_{\bar k}(M,h))$ is the poset of~\cite[Thm.~2]{AB24} described in Section~\ref{sec:partial}.
\end{theorem}

%======================================================================
\section{Local equivariant cohomology}\label{sec:local}
%======================================================================

Let $V\cong\bigoplus_{j=1}^nV(\alpha_j)$ be a representation of $T$ with all $\alpha_j\neq0$.

\begin{proposition}\label{prop:sphere}
Restriction to $0$ is an injective map $\HT(V,V\setminus0)\to\HT(0)=H^*(BT)$ with image the principal ideal $(\alpha_1\cdots\alpha_n)$.
\end{proposition}

\begin{proof}
Choose an invariant complex structure on each $V(\alpha_j)$, which orients $V$. The group $\HT(V,V\setminus0)$ is the cohomology of the Thom space of the oriented bundle $V_T=ET\times_TV\to BT$. By the Thom isomorphism~\cite{MSt74} it is free of rank one over $H^*(BT)$ on the Thom class $\tau$, and $\tau$ restricts to the Euler class $e(V_T)=\prod_j\alpha_j$ in $\HT(V)=\HT(0)$. Restriction to $0$ is therefore multiplication by $\prod_j\alpha_j$, which is injective since $H^*(BT)$ is a domain.
\end{proof}

\begin{lemma}\label{lem:ufd}
Let $k$ be a field and $R=k[t_1,\dots,t_m]$.
\begin{enumerate}[nosep,label=(\arabic*)]
\item Every nonzero linear form is irreducible, hence prime, in $R$.
\item If $\alpha_1,\dots,\alpha_n$ and $\beta_1,\dots,\beta_{n'}$ are nonzero linear forms with $(\alpha_1\cdots\alpha_n)=(\beta_1\cdots\beta_{n'})$, then $n=n'$ and there are $\sigma\in S_n$ and $c_j\in k^\times$ with $\beta_{\sigma(j)}=c_j\alpha_j$.
\end{enumerate}
\end{lemma}

\begin{proof}
(1) A factorization of a linear form has a factor of degree $0$, a unit. (2) Equal principal ideals of a domain have associate generators, and both products are factorizations into primes by~(1); uniqueness of factorization in $R$ gives the claim.
\end{proof}

\begin{theorem}\label{thm:main}
Let $pt\in X^T$ be an isolated fixed point with tangent weights $\alpha_1,\dots,\alpha_n$.
\begin{enumerate}[nosep,label=(\arabic*)]
\item The composition $\HT(X,X\setminus pt)\to\HT(X)\to\HT(pt)$ is injective, with image $(\alpha_1\cdots\alpha_n)=(e_T(T_{pt}X))$.
\item This ideal determines the multiset $\{\alpha_j\}$ up to nonzero scalars. Hence it determines the linear matroid of weights at $pt$, including parallel elements, and by Theorem~\ref{thm:AC} the lattice $S(X)_{\geqslant pt}$.
\end{enumerate}
\end{theorem}

\begin{proof}
By the slice theorem~\cite{Bre72}, $pt$ has an invariant closed neighbourhood $D$ equivariantly diffeomorphic to the unit disc of $V=T_{pt}X$. Excision of $X\setminus\operatorname{int}D$ and radial deformation give $\HT(X,X\setminus pt)\cong\HT(D,D\setminus pt)\cong\HT(V,V\setminus0)$, compatibly with restriction to $pt$. Part~(1) follows from Proposition~\ref{prop:sphere}, and part~(2) from Lemma~\ref{lem:ufd}, since rescaling does not change linear dependence.
\end{proof}

\begin{remark}
As an abstract module, $\HT(X,X\setminus pt)$ is free of rank one and carries no information; everything lies in its embedding into $\HT(pt)$. For a GKM manifold the image can also be computed from the GKM theorem, as $\bigcap_j(\alpha_j)$. Without pairwise independence this intersection differs from the product: for $\alpha_1=\alpha_2$ it is $(\alpha_1)$, not $(\alpha_1^2)$. The Thom isomorphism argument needs no such assumption, which is why Theorem~\ref{thm:main} sees multiplicities.
\end{remark}

%======================================================================
\section{Local weight matroids do not determine equivariant cohomology}\label{sec:counterexample}
%======================================================================

Let $M=\diag(1,2,3,4,5)$, $T=T^5$, and
\[
h_1=(2,3,5,5,5),\quad h_2=(2,4,4,5,5),\quad X_1=\Hess(M,h_1),\quad X_2=\Hess(M,h_2).
\]
So $X_1=\{MV_1\subset V_2,\ MV_2\subset V_3\}$ and $X_2=\{MV_1\subset V_2,\ MV_2\subset V_4,\ MV_3\subset V_4\}$. The graph $\Gamma_{h_1}$ has edges $12,23,34,35,45$: a triangle $\{3,4,5\}$ with the path $1-2-3$ attached. The graph $\Gamma_{h_2}$ has edges $12,23,24,34,45$: a triangle $\{2,3,4\}$ with two pendant edges. They are not isomorphic.

\begin{lemma}\label{lem:local-hess}
At every fixed point $w$ of $\Hess(M,h)$, the weight matroid is isomorphic to the graphic matroid $M(\Gamma_h)$.
\end{lemma}

\begin{proof}
The weights at $w$ are $e_{w(i)}-e_{w(j)}$ over the edges $\{j,i\}$ of $\Gamma_h$. The linear automorphism $e_a\mapsto e_{w^{-1}(a)}$ carries them to $e_i-e_j$, which realize $M(\Gamma_h)$.
\end{proof}

\begin{lemma}\label{lem:iso-matroids}
$M(\Gamma_{h_1})\cong M(\Gamma_{h_2})\cong U_{2,3}\oplus U_{1,1}\oplus U_{1,1}$.
\end{lemma}

\begin{proof}
In each graph the only cycle is the triangle and the other two edges are bridges.
\end{proof}

\begin{theorem}\label{thm:counterexample}
\begin{enumerate}[nosep,label=(\arabic*)]
\item $X_1$ and $X_2$ are connected smooth projective varieties of complex dimension $5$ with $H^{\mathrm{odd}}=0$. The torus $T$ acts on both with the same $120$ fixed points, as GKM actions whose effective quotient $T/\Delta(S^1)\cong T^4$ acts with complexity one.
\item For all $p_1\in X_1^T$ and $p_2\in X_2^T$ the weight matroids are isomorphic to $U_{2,3}\oplus U_{1,1}\oplus U_{1,1}$. Hence $S(X_1)_{\geqslant p_1}\cong S(X_2)_{\geqslant p_2}$, and the Chow rings and Orlik--Solomon algebras of the local matroids are isomorphic.
\item $P_{X_1}(q)=1+17q+42q^2+42q^3+17q^4+q^5$ and $P_{X_2}(q)=1+12q+47q^2+47q^3+12q^4+q^5$. In particular $\HT(X_1)\not\cong\HT(X_2)$, even as graded vector spaces.
\end{enumerate}
\end{theorem}

\begin{proof}
(3) Summing~\eqref{eq:dps} over $S_5$ gives the polynomials. Independently, the GKM presentation of $\HT(X_i)$, computed in Macaulay2 as a graded $\Q[t_1,\dots,t_5]$-module, has Hilbert series $P_{X_i}(q)/(1-q)^5$ with the same numerators. Since the actions are equivariantly formal, $\dim H^2_T(X_i)=\beta_2+5\beta_0$, which is $22$ for $X_1$ and $17$ for $X_2$.

(1) Smoothness, dimension and $H^{\mathrm{odd}}=0$ were recalled in Section~\ref{sec:prelim}; connectedness follows from $\beta_0=1$. The weights at a fixed point are distinct roots $e_a-e_b$, hence pairwise independent. An element of $T$ acting trivially lies in the common kernel of the weights $e_a-e_b$ over the edges of a connected graph, which is the diagonal circle, and the diagonal circle acts trivially on $\Fl(\C^5)$. So the effective torus has dimension $4$ and the complexity is $5-4=1$.

(2) follows from Lemmas~\ref{lem:local-hess} and~\ref{lem:iso-matroids} and Theorem~\ref{thm:AC}.
\end{proof}

\begin{remark}
A computer search over all pairs of Hessenberg functions with $n\leqslant6$ finds smaller pairs in $\Fl(\C^4)$, for example $h=(2,3,3,4)$ and $(2,2,4,4)$, but these varieties are disconnected and $\beta_0$ merely counts components. Connected pairs also exist in $\Fl(\C^6)$, such as $h=(2,3,4,6,6,6)$ and $(2,3,5,5,6,6)$, with local matroid $U_{2,3}\oplus U_{1,1}^{\oplus3}$ and Betti numbers $(1,43,179,274,179,43,1)$ and $(1,29,175,310,175,29,1)$.
\end{remark}

The global face posets of $X_1$ and $X_2$ are not isomorphic: $\Cl_{\Gamma_{h_1}}$ and $\Cl_{\Gamma_{h_2}}$ have $200$ and $190$ elements of rank~$2$. This is no accident, as the next section shows.

%======================================================================
\section{The face poset determines the Betti numbers in type A}\label{sec:rigidity}
%======================================================================

\begin{question}[\cite{Mik22}]\label{q:betti}
Let $X_1,X_2$ carry torus actions with $H^{\mathrm{odd}}(X_1)=H^{\mathrm{odd}}(X_2)=0$. Does $S(X_1)\cong S(X_2)$ imply $\beta_j(X_1)=\beta_j(X_2)$ for all $j$?
\end{question}

For quasitoric manifolds the answer is yes, since $\beta_{2i}=h_i$ of the orbit polytope. The converse fails already there: a pentagonal prism and a cube with one vertex cut off have the same $h$-vector but different face lattices.

\subsection{Rigidity of cluster-permutohedra}
For $y\in\Cl_\Gamma$ let $\nu(y)$ be the number of minimal elements below~$y$. An element $y=(C,A)$ lies above the minimal element given by a bijection $p\colon[n]\to[n]$ exactly when $p$ maps every block onto its label set, so $\nu(y)=\prod_{V_i\in C}|V_i|!$. In particular $\Cl_\Gamma$ has $n!$ minimal elements.

\begin{lemma}\label{lem:atoms-nu}
Let $x\in\Cl_\Gamma$ be minimal. Then $(\Cl_\Gamma)_{\geqslant x}\cong\Flats(\Gamma)$, and its atoms correspond to the edges of~$\Gamma$. The join of the atoms of a set $S$ of edges has rank $n-c(S)$, where $c(S)$ is the number of components of $([n],S)$. For distinct edges $e,f$, the join of their atoms has $\nu=6$ if $e$ and $f$ share a vertex and $\nu=4$ otherwise.
\end{lemma}

\begin{proof}
Above $x$ the assignment is determined by the clustering, and the join of the atoms of $S$ is the clustering into the components of $([n],S)$. For two edges this clustering has one block of size $3$, or two blocks of size $2$.
\end{proof}

\begin{theorem}\label{thm:cl-rigid}
Let $\Gamma$ and $\Gamma'$ be simple graphs on $[n]$ and $[n']$, $n,n'\geqslant1$. If $\Cl_\Gamma\cong\Cl_{\Gamma'}$, then $\Gamma\cong\Gamma'$.
\end{theorem}

\begin{proof}
Let $\Phi$ be an isomorphism. Counting minimal elements gives $n!=n'!$, so $n=n'$. Fix a minimal $x$. Then $\Phi$ restricts to an isomorphism of lattices $(\Cl_\Gamma)_{\geqslant x}\to(\Cl_{\Gamma'})_{\geqslant\Phi(x)}$ and induces a bijection $\varphi\colon E(\Gamma)\to E(\Gamma')$ on atoms. Since $\Phi$ preserves joins in these intervals, rank and $\nu$, Lemma~\ref{lem:atoms-nu} shows that $\varphi$ is an isomorphism of line graphs $L(\Gamma)\to L(\Gamma')$ preserving the rank of every set of edges. The components of $L(\Gamma)$ are the line graphs of the components $\Gamma_1,\dots,\Gamma_r$ of $\Gamma$ with at least one edge, so $\varphi$ matches them with the corresponding components $\Gamma'_i$ of $\Gamma'$ and $L(\Gamma_i)\cong L(\Gamma'_i)$. By Whitney's theorem~\cite{Whi32}, $\Gamma_i\cong\Gamma'_i$ unless one of them is $K_3$ and the other $K_{1,3}$. This is excluded: the three edges of $K_3$ span an element of rank $2$, those of $K_{1,3}$ an element of rank $3$. Both graphs have $n-\sum_i|V(\Gamma_i)|$ isolated vertices.
\end{proof}

\begin{corollary}\label{cor:q1-hess}
Let each of $X$ and $X'$ be either $M_{\Gamma,\lambda}$ with $\Gamma$ connected and $\lambda$ generic, or $\Hess(M,h)$ with $\Gamma_h$ connected. If $S(X)\cong S(X')$, then $\beta_j(X)=\beta_j(X')$ for all $j$, and $\HT(X)\cong\HT(X')$ up to an automorphism of $T$.
\end{corollary}

\begin{proof}
By Theorem~\ref{thm:AB} the face posets are $\Cl_\Gamma$ and $\Cl_{\Gamma'}$ for the corresponding graphs, so $\Gamma\cong\Gamma'$ by Theorem~\ref{thm:cl-rigid}. If $\pi$ is a permutation with $\pi(\Gamma)=\Gamma'$ and $P_\pi$ its permutation matrix, conjugation $A\mapsto P_\pi AP_\pi^{-1}$ preserves the spectrum and the zero pattern up to relabelling, $(P_\pi AP_\pi^{-1})_{\pi(i)\pi(j)}=a_{ij}$. It is therefore a diffeomorphism $M_{\Gamma,\lambda}\to M_{\Gamma',\lambda}$, equivariant along the automorphism $t\mapsto P_\pi tP_\pi^{-1}$ of $T$. A Hessenberg variety and its twin $M_{\Gamma_h,\lambda}$ have isomorphic equivariant cohomology and equal Betti numbers~\cite[Thm.~3.11]{AB21}.
\end{proof}

The corollary does not use $H^{\mathrm{odd}}=0$. In particular it covers $M_{K_{1,3},\lambda}$ and the manifolds $M_{C_n,\lambda}$, $n\geqslant4$, which are not equivariantly formal~\cite{AB24,AMS23}.

\subsection{Partial Hessenberg varieties}\label{sec:partial}
Let $X=\Hess_{\bar k}(M,h)$. By~\cite[Thm.~2]{AB24}, its faces are the sets $\{(W_j,c_j)\}$ in which the blocks $W_j$ partition $[n]$, each $c_j$ is a vector of part counts with $\sum_ac_j(a)=|W_j|$ and $\sum_jc_j=\bar k$, and each block induces a connected subgraph of the inflation. The order is refinement with counts adding up, and the rank is $n$ minus the number of blocks. The minimal faces are the fixed points, and the number of fixed points below a face is
\[
\nu\bigl(\{(W_j,c_j)\}\bigr)=\prod_j\frac{|W_j|!}{\prod_ac_j(a)!}.
\]
For every fixed point $J$, $S(X)_{\geqslant J}\cong\Flats(\Gamma_h(\bar k))$, in accordance with Theorem~\ref{thm:AC} and the weights of Section~\ref{sec:prelim}.

\begin{theorem}\label{thm:q1-partial}
If $X$ and $X'$ are partial regular semisimple Hessenberg varieties with $S(X)\cong S(X')$, then $\beta_j(X)=\beta_j(X')$ for all~$j$.
\end{theorem}

The proof rests on three lemmas.

\begin{lemma}[Splitting]\label{lem:split}
Let $Q_1,\dots,Q_r$ be the connected components of $\Gamma_h$ and $P_{Q_i}$ the Poincaré polynomial of the partial Hessenberg variety of the restriction of $(\bar k,h)$ to~$Q_i$. Then
\[
P_X(q)=\#X^T\cdot\prod_{i=1}^r\frac{P_{Q_i}(q)}{P_{Q_i}(1)},
\]
and a component consisting of a single part contributes the factor~$1$.
\end{lemma}

\begin{proof}
The components of the indifference graph $\Gamma_h$ are intervals of parts, and every pair counted in~\eqref{eq:bb} lies in one component. Put $U_i=\bigcup_{a\in Q_i}J_a$. The pairs counted inside $Q_i$ depend only on the relative order of the elements of $U_i$, that is, on a fixed point of the variety of $Q_i$. Summing over the choices of $U_1,\dots,U_r$ gives $P_X=\binom{n}{|U_1|,\dots,|U_r|}\prod_iP_{Q_i}$, and setting $q=1$ gives $\#X^T=\binom{n}{|U_1|,\dots,|U_r|}\prod_iP_{Q_i}(1)$.
\end{proof}

\begin{lemma}[Fibration]\label{lem:fibration}
Put $\tilde h(i)=K_{h(a)}$ for $K_{a-1}<i\leqslant K_a$, and let $\pi\colon\Fl(\C^n)\to\Fl_{\bar k}$ be the projection. Then $\pi^{-1}(X)=\Hess(M,\tilde h)$, the restriction of $\pi$ is a fibre bundle with fibre $\prod_a\Fl(\C^{k_a})$, and
\[
P_{\Hess(M,\tilde h)}(q)=P_X(q)\cdot\prod_{a=1}^s[k_a]_q!.
\]
Moreover, $\Gamma_{\tilde h}$ is the inflation $\Gamma_h(\bar k)$ with every part made a clique.
\end{lemma}

\begin{proof}
A full flag $F_\bullet$ lies in $\Hess(M,\tilde h)$ iff $MF_i\subset F_{\tilde h(i)}$ for all~$i$. For $i$ in part $a$ we have $F_i\subset F_{K_a}$ and $\tilde h(i)=\tilde h(K_a)$, so the condition at $i$ follows from the condition at $K_a$, and the conditions at $K_1,\dots,K_s$ say that $\pi(F_\bullet)\in X$. The projection $\pi$ is a holomorphic bundle with fibre $\prod_a\Fl(V_{K_a}/V_{K_{a-1}})$ over $V_\bullet$, so it restricts to a bundle over~$X$. The cohomology of the fibre is generated by Chern classes of tautological bundles defined on the total space, and all spaces involved have only even cohomology, so the Leray--Hirsch theorem gives the product formula, with $P_{\Fl(\C^k)}=[k]_q!$. Finally, $i<j$ is an edge of $\Gamma_{\tilde h}$ iff $j\leqslant K_{h(a)}$ for the part $a$ of $i$, that is, iff $j$ lies in part $a$ or in a part $b$ with $a<b\leqslant h(a)$.
\end{proof}

\begin{lemma}[Reconstruction]\label{lem:reconstruct}
Let $G=\Gamma_h(\bar k)$ with its partition into parts. The poset $S(X)$ determines $\#X^T$ and the multiset of isomorphism classes, as graphs with a vertex partition, of the components of $G$ that have an edge.
\end{lemma}

\begin{proof}
The number $\#X^T$ is the number of minimal elements. Fix a fixed point $J$ and identify $S(X)_{\geqslant J}$ with $\Flats(G)$. For edges $uv$ and $vw$ the join of their atoms has the single block $\{u,v,w\}$, so $\nu=3$ if $u$ and $w$ lie in one part and $\nu=6$ otherwise, since adjacent vertices lie in different parts. Two disjoint edges give $\nu=4$. As in the proof of Theorem~\ref{thm:cl-rigid}, an isomorphism $\Phi\colon S(X)\to S(X')$ induces a bijection $\varphi$ between the edge sets of the inflations that preserves the rank and $\nu$ of the join of every set of edges. In particular $\varphi$ is an isomorphism of line graphs. The components of $G$ with an edge are the inflations $G_Q$ of the components $Q$ of $\Gamma_h$ with at least two parts, and copies of a part that is a component of $\Gamma_h$ by itself are isolated. So $\varphi$ matches these components. Let $G_Q$ and $G'_{Q'}$ be matched.

If $G_Q$ has more than four vertices, then so has $G'_{Q'}\cong G_Q$, and by Jung's theorem~\cite{Jun66} $\varphi|_{E(G_Q)}$ is induced by a graph isomorphism $\psi$. Two distinct vertices $u,w$ of $G_Q$ lie in one part iff they have a common neighbour $v$ with $\nu(a_{uv}\vee a_{vw})=3$; copies of one part do have a common neighbour, as $Q$ is connected with at least two parts. Since $\psi$ preserves adjacency and $\varphi$ preserves $\nu$, $\psi$ maps parts to parts.

If $G_Q$ has at most four vertices, then so has $G'_{Q'}$, since by Whitney's theorem they are isomorphic or are $K_3$ and $K_{1,3}$. Each is the inflation of a connected Hessenberg graph with $n\leqslant4$. A computer check over all of them ($13$ up to isomorphism of graphs with parts) shows that an edge bijection between two of them that preserves rank and $\nu$ of the joins of all edge sets exists only if they are isomorphic as graphs with parts.
\end{proof}

\begin{proof}[Proof of Theorem~\ref{thm:q1-partial}]
By Lemma~\ref{lem:split} it suffices to show that $S(X)$ determines $\#X^T$ and $\prod_QP_Q(q)/P_Q(1)$ over the components $Q$ of $\Gamma_h$ with at least two parts. For such $Q$, Lemma~\ref{lem:fibration} gives $P_Q=P_{\Hess(M_Q,\tilde h_Q)}/\prod_{a\in Q}[k_a]_q!$. The graph $\Gamma_{\tilde h_Q}$ is connected and is obtained from $G_Q$ by making the parts cliques, so by Corollary~\ref{cor:q1-hess} and Theorem~\ref{thm:AB} its Poincaré polynomial depends only on $G_Q$ with its parts, up to isomorphism; so do the part sizes. Lemma~\ref{lem:reconstruct} completes the proof.
\end{proof}

\subsection{Remarks}
Theorem~\ref{thm:q1-partial} includes partial flag varieties ($\Gamma_h=K_s$) and Grassmannians, where it is classical, and the full case $\bar k=(1,\dots,1)$. For disconnected $\Gamma_h$ the variety $\Hess(M,h)$ is a disjoint union of products of smaller Hessenberg varieties, and Theorem~\ref{thm:q1-partial} covers it as well.

Lemma~\ref{lem:fibration} and Theorem~\ref{thm:cl-rigid} were also checked by computer: the identity of Lemma~\ref{lem:fibration} for all $2\,949$ pairs $(\bar k,h)$ with $n\leqslant7$, and Theorem~\ref{thm:cl-rigid} by exact poset isomorphism for all $208$ graphs with $n\leqslant6$.

%======================================================================
\section{Open questions}\label{sec:open}
%======================================================================

\begin{enumerate}[label=(\arabic*)]
\item \emph{Other Lie types.} Regular semisimple Hessenberg varieties exist for every reductive group. Is there an analogue of Theorem~\ref{thm:cl-rigid} for their face posets? Outside type~A the local weight matroids are no longer graphic, and Whitney's theorem is not available.
\item \emph{A formula.} In complexity zero, $S(X)$ determines the Betti numbers through the $h$-vector. Is there a formula for $\beta_j$ in terms of $S(X)$ for Hessenberg varieties, as opposed to the reconstruction argument given here?
\item \emph{Beyond algebraic actions.} Question~\ref{q:betti} remains open for general equivariantly formal torus manifolds of positive complexity.
\item \emph{Oriented weights.} If $X$ carries an invariant stably complex structure, the weights have signs and $S(X)_{\geqslant p}$ refines to an oriented matroid. What does the resulting oriented structure determine?
\end{enumerate}

\paragraph{Acknowledgements.}
The author is grateful to V.~M.~Buchstaber, the advisor of the thesis~\cite{Mik22} from which this work grew.
% TODO(author): further acknowledgements, affiliation and e-mail.

\begin{thebibliography}{99}\small

\bibitem{AH20} H.~Abe, T.~Horiguchi, \emph{A survey of recent developments on Hessenberg varieties}, in: Schubert Calculus and Its Applications in Combinatorics and Representation Theory, Springer Proc.\ Math.\ Stat., Springer, Singapore, 2020, 251--279.

\bibitem{AHK18} K.~Adiprasito, J.~Huh, E.~Katz, \emph{Hodge theory for combinatorial geometries}, Ann.\ of Math.\ (2) 188 (2018), no.~2, 381--452.

\bibitem{AB21} A.~Ayzenberg, V.~Buchstaber, \emph{Manifolds of isospectral matrices and Hessenberg varieties}, Int.\ Math.\ Res.\ Not.\ IMRN 2021, no.~21, 16671--16692.

\bibitem{AB24} A.~Ayzenberg, V.~Buchstaber, \emph{Cluster-permutohedra and submanifolds of flag varieties with torus actions}, Int.\ Math.\ Res.\ Not.\ IMRN 2024, no.~3, 1931--1967.

\bibitem{AC22} A.~Ayzenberg, V.~Cherepanov, \emph{Matroids in toric topology}, arXiv:2203.06282.

\bibitem{AMS23} A.~Ayzenberg, M.~Masuda, G.~Solomadin, \emph{How is a graph not like a manifold?}, Sb.\ Math.\ 214 (2023), no.~6, 793--815.

\bibitem{Bre72} G.~E.~Bredon, \emph{Introduction to Compact Transformation Groups}, Pure Appl.\ Math.\ 46, Academic Press, 1972.

\bibitem{BP15} V.~M.~Buchstaber, T.~E.~Panov, \emph{Toric Topology}, Math.\ Surveys Monogr.\ 204, Amer.\ Math.\ Soc., 2015.

\bibitem{BT19a} V.~M.~Buchstaber, S.~Terzić, \emph{Toric topology of the complex Grassmann manifolds}, Mosc.\ Math.\ J.\ 19 (2019), no.~3, 397--463.

\bibitem{BT19b} V.~M.~Buchstaber, S.~Terzić, \emph{Foundations of $(2n,k)$-manifolds}, Mat.\ Sb.\ 210 (2019), no.~4, 41--86 (in Russian); English transl.\ in Sb.\ Math.\ 210 (2019), no.~4.

\bibitem{CS74} T.~Chang, T.~Skjelbred, \emph{The topological Schur lemma and related results}, Ann.\ of Math.\ (2) 100 (1974), 307--321.

\bibitem{DPS92} F.~De~Mari, C.~Procesi, M.~A.~Shayman, \emph{Hessenberg varieties}, Trans.\ Amer.\ Math.\ Soc.\ 332 (1992), no.~2, 529--534.

\bibitem{GKM98} M.~Goresky, R.~Kottwitz, R.~MacPherson, \emph{Equivariant cohomology, Koszul duality, and the localization theorem}, Invent.\ Math.\ 131 (1998), no.~1, 25--83.

\bibitem{M2} D.~R.~Grayson, M.~E.~Stillman, \emph{Macaulay2, a software system for research in algebraic geometry}, available at \url{https://macaulay2.com}.

\bibitem{Jun66} H.~A.~Jung, \emph{Zu einem Isomorphiesatz von H.~Whitney für Graphen}, Math.\ Ann.\ 164 (1966), 270--271.

\bibitem{Mik22} D.~A.~Mikheev, \emph{Invariants of matroids defined by torus actions}, diploma thesis, Lomonosov Moscow State University, 2022 (in Russian); revised English edition, 2026, \url{https://github.com/davidmiheev/matroid_toric}.

\bibitem{MP14} B.~D.~McKay, A.~Piperno, \emph{Practical graph isomorphism, II}, J.\ Symbolic Comput.\ 60 (2014), 94--112.

\bibitem{MSt74} J.~W.~Milnor, J.~D.~Stasheff, \emph{Characteristic Classes}, Ann.\ of Math.\ Stud.\ 76, Princeton Univ.\ Press, 1974.

\bibitem{OT92} P.~Orlik, H.~Terao, \emph{Arrangements of Hyperplanes}, Grundlehren Math.\ Wiss.\ 300, Springer, 1992.

\bibitem{Tym08} J.~S.~Tymoczko, \emph{Permutation actions on equivariant cohomology of flag varieties}, in: Toric Topology, Contemp.\ Math.\ 460, Amer.\ Math.\ Soc., 2008, 365--384.

\bibitem{Whi32} H.~Whitney, \emph{Congruent graphs and the connectivity of graphs}, Amer.\ J.\ Math.\ 54 (1932), no.~1, 150--168.

\end{thebibliography}
\end{document}
